% !TeX root = ../../Book2.tex
\chapter{Fourier 乘子}
\label{b2:ch:22}
\BookChapterNumber{b2:ch:22}
\begin{chapterabstract}
本章把 Fourier 侧的有界符号转化为 \(L^p\) 算子。完整证明一个有限导数版本的 Mihlin 定理及 Sobolev 环带版本的 Hörmander 定理，并应用于 Riesz 变换、Bessel 位势、平方函数正则性和常系数椭圆估计。
\end{chapterabstract}

\begin{learninggoals}\begin{itemize}
\item\label{b2:item:22-goal-1}掌握 \(L^2\) 乘子的准确刻画及尺度、对偶和复合规则。
\item\label{b2:item:22-goal-2}能把尺度导数条件转换成频率环带核的统一估计。
\item\label{b2:item:22-goal-3}理解 Hörmander 积分差分条件如何取代核的逐点光滑条件。
\item\label{b2:item:22-goal-4}能区分同指数乘子界、分数积分界和带导数的椭圆先验估计。
\end{itemize}\end{learninggoals}

卷积算子在 Fourier 一侧成为乘法，但 \(L^p\) 范数并不被 Fourier 变换保持。\chapref{b2:ch:18}建立的奇异积分理论与\chapref{b2:ch:19}的频率分解共同弥补这一缺口：每个环带有快速衰减核，所有尺度的叠加则通过零阶奇异积分处理。以下定理中的导数数目是明确的充分条件，凡未证明的锐利改进都不纳入正文前置。

% !TeX root = ../../Book2.tex
\section{乘子算子}
\label{b2:sec:2201}

平移不变算子在 Fourier 一侧表现为逐点乘法。本章讨论怎样从符号的局部光滑性得到整个 \(L^p\) 范数界，并区别阶数为零的有界乘子、改变可积指数的分数积分和具有导数损失的微分算子。

\subsection{Fourier 乘子}
\label{b2:subsec:220101}

\begin{definition}\label{b2:def:22-multiplier}
对可测函数 \(m\)，先在使乘积有意义的测试函数上定义
\[
 T_m f=\mathcal F^{-1}(m\widehat f).
\]
若该算子延拓为 \(L^p(\mathbb R^d)\) 上的有界算子，称 \(m\) 为 \(L^p\) Fourier 乘子，并记其算子范数为 \(\|m\|_{M_p}\)。
\end{definition}
乘子在零测集上的值不影响算子。对于一般有多项式增长的光滑 \(m\)，乘法可在 Schwartz 空间与温和分布之间解释；对于仅可测的有界 \(m\)，最稳妥的起点是 \(L^2\)。
若 \(m=\widehat k\)、\(k\in L^1\)，则 \(T_mf=k*f\)，Young 不等式给出所有 \(1\le p\le\infty\) 上的有界性。

\ProseParagraphBreak
乘法和卷积的对应不意味着每个乘子都有普通可积核。恒等乘子的核为 \(\delta_0\)，Hilbert 乘子的核为主值分布。推导核估计时应先局部化频率，再在原点外重构普通函数；原点上的分布项要由已有的 \(L^2\) 定义处理。

\subsection{\texorpdfstring{\(L^2\)}{L2} 理论}
\label{b2:subsec:220102}

\begin{proposition}\label{b2:prop:22-l2}
可测 \(m\) 是 \(L^2\) 乘子当且仅当 \(m\in L^\infty\)，且
\(\|m\|_{M_2}=\|m\|_\infty\)。
\end{proposition}
\begin{proof}
Plancherel 给出充分性及范数上界。若算子范数为 \(C\)，在任意有限测度集合 \(E\) 上选 \(\widehat f=\mathbf1_E\)，得
\(\int_E|m|^2\le C^2|E|\)。若 \(|m|>C+\epsilon\) 的集合有正测度，取其中有限测度的正测度子集就矛盾。因此 \(|m|\le C\) 几乎处处，得到必要性和等号。
\end{proof}
对 \(|m|=1\) 的符号，算子在 \(L^2\) 上等距；这并不提供其他 \(p\) 上的范数控制。即使每个频率只发生相位旋转，不同频率叠加后的空间集中仍可能发生变化。

\ProseParagraphBreak
若 \(m_1,m_2\in M_p\)，则 \(m_1m_2\in M_p\)，且范数不超过两范数之积。先在 \(L^2\cap L^p\) 上验证 \(T_{m_1m_2}=T_{m_1}T_{m_2}\)，再用稠密性延拓即可。相似地，乘子空间对一致算子范数有界的点态极限封闭：在 \(L^2\cap L^p\) 上先取 \(L^2\) 极限及几乎处处子列，再用 Fatou，最后延拓。这是处理光滑截断极限的常用办法。

\subsection{\texorpdfstring{\(L^p\)}{Lp} 乘子问题}
\label{b2:subsec:220103}

若 \(m\in M_p\)、\(1<p<\infty\)，其 Banach 伴随在 \(L^{p'}\) 上的符号为 \(\overline m\)，范数相同。若 \(m\) 实值，就得到对偶指数上的相同范数。与 \(L^2\) 界插值，可控制 \(p\) 和 \(2\) 之间的全部指数，但不能越过尚未建立的端点。

\ProseParagraphBreak
缩放不改变乘子范数。若 \(D_rf(x)=f(rx)\)，则
\(T_{m(r\,\cdot)}=D_{1/r}T_mD_r\)，而两个缩放范数相消。
正交变换与频率平移也保持范数：前者对应空间坐标旋转，后者对应调制。符号乘上 \(e^{-2\pi\mathrm i y\cdot\xi}\) 则对应输出平移。这些等距对称使乘子条件应具有尺度不变性；要求所有导数在整个频率空间绝对有界，并不是适合齐次问题的条件。

\ProseParagraphBreak
一个不光滑却有界的例子是一维区间特征函数。半轴投影可用恒等算子与 Hilbert 变换表示，两个频率平移后的半轴投影之差就是区间投影，因此在 \(1<p<\infty\) 上一致有界。由此不能推断高维球的特征函数同样有界。符号边界的几何维数与曲率会参与问题，\chapref{b2:ch:24}将专门说明这一障碍。

% 扩写定位：b2:subsec:220103
Fourier 乘子并不是凭空选出的一类算子；在 \(L^2\) 上，它们正好是与所有平移交换的有界线性算子。
\begin{proposition}\label{b2:prop:22-translation-multiplier}
若 \(T:L^2(\mathbb R^d)\to L^2(\mathbb R^d)\) 有界，且 \(T\tau_y=\tau_yT\) 对全部 \(y\) 成立，其中 \(\tau_yf(x)=f(x-y)\)，则存在唯一的 \(m\in L^\infty\)，使 \(T=T_m\)，并且 \(\|T\|_{2\to2}=\|m\|_\infty\)。
\end{proposition}
\begin{proof}
令 \(A=\mathcal FT\mathcal F^{-1}\)。平移变成乘法 \(e^{-2\pi\mathrm iy\cdot\xi}\)，故 \(A\) 与所有这些乘法算子交换。对任意 \(\varphi\in C_c^\infty\)，Fourier 反演把 \(\varphi\) 写成这些指数乘法的 \(L^1\) 加权积分；其在 \(L^2\) 上的强积分与 \(A\) 交换。于是 \(A\) 与乘法 \(M_\varphi\) 交换。

由光滑函数的有界几乎处逼近，\(A\) 也与任意可测集合的投影 \(M_{\mathbf1_E}\) 交换。可以先对有界集合用 Lebesgue 正则性与卷积逼近得到 \(0\leq\varphi_n\leq1\)、\(\varphi_n\to\mathbf1_E\) 几乎处，支配收敛给出强算子收敛；再以球截断推广到一般集合。

取递增的有限测度立方体 \(Q_n\) 覆盖全空间。在 \(Q_n\) 上定义 \(m_n=A\mathbf1_{Q_n}\)。若 \(n<k\)，交换投影给出
\(\mathbf1_{Q_n}m_k=A(\mathbf1_{Q_n}\mathbf1_{Q_k})=m_n\)，所以拼成可测函数 \(m\)。对 \(E\subset Q_n\)，有
\[
 A\mathbf1_E=\mathbf1_E m,\qquad
 \int_E|m|^2\leq\|A\|^2|E|.
\]
若 \(\{|m|>\|A\|+\epsilon\}\cap Q_n\) 有正测度，将其代入即矛盾。因此 \(\|m\|_\infty\leq\|A\|\)。有限支集简单函数的稠密性说明 \(A=M_m\)，反向范数与唯一性由本节的 \(L^2\) 乘子结论得到。
\end{proof}
该结论只建立在 \(L^2\) 上。若算子只在 \(L^p\) 给定，还不能通过一次 Fourier 变换就把它视为 \(L^2\) 有界乘子；需先说明它在共同稠密类上的作用及其额外的平方界。另一方面，具有不连续符号的平移不变算子在 \(L^2\) 上完全没有困难；真正的乘子问题是同一个算子能否保持其他 \(L^p\)。

% !TeX root = ../../Book2.tex
\section{Mihlin–Hörmander 定理}
\label{b2:sec:2202}

导数条件的作用是使每个频率环带上的核有足够快的衰减，再让环带总和满足奇异积分的抵消估计。证明须同时使用符号的有界性来建立 \(L^2\) 起点；核的大小条件本身不能替代这个起点。

\subsection{符号导数}
\label{b2:subsec:220201}

设 \(N\) 为整数。常用的尺度不变符号界为
\[
 |\partial^\alpha m(\xi)|\le A|\xi|^{-|\alpha|},
 \qquad 0\le|\alpha|\le N,\quad \xi\ne0.
\]
零阶包含 \(|m|\le A\)。例如 \(\xi_j/|\xi|\) 的各阶导数都是相应次数的齐次函数，在单位球面连续有界，因此满足任意有限阶条件。相反，\(|\xi|^{-\alpha}\) 在零频无界，不能成为 \(L^2\) 上的有界零阶乘子。

\ProseParagraphBreak
若 \(m(\xi)=b(\log|\xi|)\)，而 \(b\) 的有限阶导数有界，则链式法则给出同样的尺度界。符号不需要在零点有极限，甚至可以沿尺度持续振荡。这里对零点的忽略不影响 \(L^p\) 算子，却提醒我们不能把定理误述成“符号在闭空间上光滑”。

\subsection{二进分解}
\label{b2:subsec:220202}

使用\chapref{b2:ch:19}的环带分割 \(\varphi(2^{-j}\xi)\)，令
\[
 m_j(\xi)=m(\xi)\varphi(2^{-j}\xi),\qquad
 k_j(x)=2^{jd}h_j(2^jx),\quad
 h_j=\mathcal F^{-1}\bigl[m(2^j\,\cdot)\varphi\bigr].
\]
在固定环带上，归一化符号的有限阶导数一致有界。对 Fourier 积分分部积分，并分别选择绝对值最大的坐标，可得
\[
 |h_j(x)|+|\nabla h_j(x)|\le C A(1+|x|)^{-N}.
\]
梯度估计只是把符号乘上有界环带中的 \(2\pi\mathrm i\xi_\ell\)，仍使用相同阶数的符号导数；不需要额外的一阶 \(m\) 导数。该估计对所有 \(j\) 一致。

\ProseParagraphBreak
当 \(N>d+1\)，在 \(x\ne0\) 上绝对求和：
\[
 \sum_j|k_j(x)|\le CA|x|^{-d},\qquad
 \sum_j|\nabla k_j(x)|\le CA|x|^{-d-1}.
\]
例如第一式把 \(2^j|x|\le1\) 与 \(>1\) 分开，分别求等比级数
\(\sum2^{jd}\) 和 \(|x|^{-N}\sum2^{j(d-N)}\)。
第二式同理需要 \(N>d+1\)。这些级数在离原点的紧集上一致收敛，从而其和 \(K\) 是那里的一阶光滑函数。

\subsection{Mihlin 乘子定理}
\label{b2:subsec:220203}

下面采用\term[mihlin-chengzi-dingli]{Mihlin 乘子定理}[Mihlin multiplier theorem]的下述版本。

\begin{theorem}[Mihlin 乘子定理的有限导数版本]\label{b2:thm:22-mihlin}
设 \(m\in C^{d+2}(\mathbb R^d\setminus\{0\})\)，且对
\(|\alpha|\le d+2\) 满足上述尺度导数界。则
\[
 \|T_mf\|_p\le C_{d,p}A\|f\|_p,\qquad1<p<\infty.
\]
该导数数目是本证明使用的方便充分条件，并不声称最优。
\end{theorem}
\begin{proof}
Plancherel 给出 \(L^2\) 范数至多 \(A\)。有限环带和的符号在几乎处处频率上趋于 \(m\)，且一致有界，因此对应算子在 \(L^2\) 中趋于 \(T_m\)。
取紧支撑 \(f\in L^2\)，其支撑之外的任意紧集到支撑的距离为正。前述核级数在那里一致绝对可积，可交换求和和积分，得到
\(T_mf(x)=\int K(x-y)f(y)\,\mathrm dy\)。
更准确地，先在两个不交支撑的测试函数之间配对，再取 \(L^2\) 极限，即确认这就是给定算子的支撑外核，而不是另行定义的算子。

核的大小与梯度界满足\chapref{b2:ch:18}的标准 Calderón--Zygmund 条件，且卷积核的两变量差分由同一个梯度界控制。\cref{b2:thm:18-cz-lp}因而给出所有 \(1<p<\infty\) 上的界。对 \(L^2\cap L^p\) 的界先建立，再以该空间在 \(L^p\) 中的稠密性定义唯一延拓。
\end{proof}
证明实际上还给出弱型 \((1,1)\)，但通常没有强型 \((1,1)\) 或 \(L^\infty\to L^\infty\) 界。Hilbert 变换的区间输入产生对数尾和边界对数奇性，已分别说明这两个端点为何不能由定理自动补齐。

\ProseParagraphBreak
有限频率环带上的导数条件也可使用：若 \(m\) 支撑在固定环带并有 \(d+1\) 阶导数，则它的核属于 \(L^1\)，从而包括端点在内都可有界。真正造成端点困难的是无限尺度叠加，不能把单个环带的 Young 估计在没有可求和常数时直接相加。

\subsection{Hörmander 条件}
\label{b2:subsec:220204}

导数的逐点控制可改为环带上的 Sobolev 控制。固定 \(\eta\in C_c^\infty\)，在 \(\varphi\) 支撑邻域等于 \(1\)，令
\[
 B_s(m)=\sup_{j\in\mathbb Z}\|m(2^j\,\cdot)\eta\|_{H^s(\mathbb R^d)}.
\]
这里 \(H^s\) 的变量是频率变量本身；此处的 Sobolev 范数不是输入 \(f\) 的范数。
下面采用\term[hormander-chengzi-dingli]{Hörmander 乘子定理}[Hörmander multiplier theorem]的下述版本。

\begin{theorem}[Hörmander 环带条件]\label{b2:thm:22-hormander}
若 \(s>d/2\)、\(m\in L^\infty\) 且 \(B_s(m)<\infty\)，则 \(m\in M_p\)、\(1<p<\infty\)，算子范数不超过
\(C_{d,p,s}(\|m\|_\infty+B_s(m))\)。
\end{theorem}
\begin{proof}
设 \(h_j\) 如上一小节。Sobolev 范数经 Plancherel 等价于
\(\|(1+|x|)^s h_j\|_2\)；固定光滑截断的乘法在 \(H^s\) 有界，故这些范数一致受 \(CB_s\) 控制。
证明这一乘法事实，只需把乘积的 Fourier 变换写为卷积，使用
\((1+|x|)^s\le C_s(1+|x-y|)^s(1+|y|)^s\) 及 Schwartz 核的带权 \(L^1\) 范数。
同样，乘上环带中的坐标 \(\xi_\ell\) 表明
\(\|(1+|x|)^s\nabla h_j\|_2\le CB_s\)。
由 Cauchy–Schwarz，
\[
 \|h_j\|_1+\|\nabla h_j\|_1\le CB_s,\qquad
 \int_{|x|>R}|h_j(x)|\,\mathrm dx\le CB_sR^{-(s-d/2)}
 \quad(R\ge1).
\]

记 \(\gamma=s-d/2>0\)。对位移 \(v\ne0\)，低频环带 \(2^j|v|\le1\) 满足
\(\|k_j(\,\cdot-v)-k_j\|_1\le CB_s2^j|v|\)。
高频环带 \(2^j|v|>1\) 在 \(|x|>2|v|\) 上用两个尾积分，得到
\[
 \int_{|x|>2|v|}|k_j(x-v)-k_j(x)|\,\mathrm dx
 \le CB_s(2^j|v|)^{-\gamma}.
\]
两段等比级数可求和，故每个有限环带和的核及其伴随都满足一致的积分 Hörmander 条件。

\chapref{b2:ch:18}弱型证明中，坏块的估计只用该积分条件：用零平均把核减去块中心的值，再把差分积分界乘以块的 \(L^1\) 范数。好部仍用一致的 \(L^2\) 范数。因此同一证明给出有限和的一致弱型界；Marcinkiewicz 插值和伴随对偶给出一致的 \(L^p\) 界。
有限和在 \(L^2\) 中趋于 \(T_m\)，取几乎处处子列并用 Fatou，便将 \(L^p\) 界传到极限，最后用稠密性延拓。
这条证明不要求极限核有逐点梯度界，因而确实比上一版本弱。
\end{proof}
严格条件 \(s>d/2\) 用于尾积分中的正指数 \(\gamma\)。在等号处，高频环带级数不再衰减；这说明本证明失效的位置，不能把失效本身当作等号处任何可能强化的反例。

% 扩写定位：b2:subsec:220204
有限阶条件中的尺度因子不可删除。若在频率 \(\xi\) 附近的变化尺度是 \(|\xi|\)，则一次求导自然带来 \(|\xi|^{-1}\)；Mihlin 假设以 \(|\xi|^{|\alpha|}\) 恢复这个尺度。普通的“全部导数有界”在无穷频率未必提供这种衰减，也就不能直接控制重缩放后的二进核。

\ProseParagraphBreak

例如 \(m(\xi)=e^{2\pi\mathrm ia\cdot\xi}\) 的乘子是平移，故在每个 \(L^p\) 上都是等距。但其导数在无穷处不衰减，除非 \(a=0\)，不满足 Mihlin 的尺度导数条件。这个例子表明定理提供的是充分条件，而不是所有有界乘子的分类。反过来，把一个任意有界但高度不规则的符号放入 \(L^\infty\)，只能保证 \(L^2\) 有界，不能得到其他指数。

\ProseParagraphBreak

还应区分单个乘子和一族乘子的界。若 \(m_\nu\) 的 Mihlin 常数一致，则标量算子范数一致；若要在位置上先对输出求平方和，需要更强的向量论证。不过在此处统一符号条件已经足以给出这一论证：对有限族取对角乘子 \(\operatorname{diag}(m_\nu)\)，归一化后的核的算子范数由各个标量核上界的上确界控制，\(L^2(\ell^2)\) 界由逐项 Plancherel 相加得到。\cref{b2:prop:18-hilbert-valued-cz}于是给出
\[
 \left\|\left(\sum_\nu|T_{m_\nu}f_\nu|^2\right)^{1/2}\right\|_p
 \leq C_p\left\|\left(\sum_\nu|f_\nu|^2\right)^{1/2}\right\|_p.
\]
该常数与族的大小无关。若只知道各个算子的 \(L^p\) 范数一致，而不知道统一核或其他向量结构，则不能凭空推出上述平方和界。这也是前面 Bessel 平方函数识别时应核对的额外条件。

% !TeX root = ../../Book2.tex
\section{应用}
\label{b2:sec:2203}

应用乘子定理时，先判断符号阶数，再选取匹配的函数空间。零阶符号可保持 \(L^p\)，负阶符号在不同频率区间可能改善正则性或改变可积指数，正阶符号则需输入具有相应导数。

\subsection{Riesz 变换}
\label{b2:subsec:220301}

Riesz 变换的符号 \(-\mathrm i\xi_j/|\xi|\) 满足任意阶的尺度导数条件，故\chapref{b2:ch:18}的 \(L^p\) 定理也可由 Mihlin 定理再次得到。两条路线揭示不同方面：奇异积分证明直接看到零平均和核尾部；乘子证明利用频率光滑性，以 \(L^2\) 自动建立算子起点。

\ProseParagraphBreak
还可处理任意光滑零次齐次符号 \(m\)，因为单位球面上的有限阶导数有界。零次齐次性并不要求球面平均为零；常数符号的核有一个原点质量。只有在提取原点外的奇异核时，才需要辨认其球面抵消，不能把抵消条件错误地加到所有零阶符号上。

\subsection{分数幂}
\label{b2:subsec:220302}

记
\[
 J^sf=\mathcal F^{-1}\bigl[(1+4\pi^2|\xi|^2)^{s/2}\widehat f\bigr],
 \qquad (-\Delta)^{s/2}f
 =\mathcal F^{-1}\bigl[(2\pi|\xi|)^s\widehat f\bigr].
\]
前者是不齐次分数阶算子，后者是齐次分数阶算子。对 \(s>0\)，\(J^{-s}\) 的符号满足 Mihlin 条件，故在 \(L^p\) 上有界；\((-\Delta)^{-s/2}\) 在零频发散，应在适当齐次空间或以\chapref{b2:ch:19}的分数积分解释，不能据此声称它在同一个 \(L^p\) 上有界。

\ProseParagraphBreak
对 \(1<p<\infty\)，定义 Bessel 位势空间
\[
 H^{s,p}=\{f\in\mathcal S':J^sf\in L^p\},\qquad
 \|f\|_{H^{s,p}}=\|J^sf\|_p.
\]
这是完整空间，因为 \(J^{-s}\) 把 \(L^p\) 连续送入温和分布，\(L^p\) 极限可经它反解。连续性的核对可对任意测试函数 \(\psi\) 使用
\(|\langle J^{-s}g,\psi\rangle|\le\|g\|_p\|J^{-s}\psi\|_{p'}\)；
任意实 \(s\) 时 \(J^{-s}\psi\) 仍为 Schwartz 函数。

\begin{proposition}\label{b2:prop:22-bessel-square}
对非齐次分割 \(P_0,\Delta_j\)（\(j\ge1\)），有
\[
 \|J^sf\|_p\asymp
 \left\|\left(|P_0f|^2+\sum_{j\ge1}2^{2js}|\Delta_jf|^2\right)^{1/2}\right\|_p.
\]
因而\chapref{b2:ch:19}定义的 \(F^s_{p,2}\) 与 \(H^{s,p}\) 相同，范数等价。
\end{proposition}
\begin{proof}
设 \(g=J^sf\)。第 \(j\) 个加权块的符号为
\(2^{js}\varphi(2^{-j}\xi)(1+4\pi^2|\xi|^2)^{-s/2}\)。
在归一化环带上，它及有限阶导数一致有界，核具有\chapref{b2:ch:19}向量值平方函数证明使用的一致大小、光滑与 \(L^2(\ell^2)\) 界。于是该证明给出
\(\|(\sum2^{2js}|\Delta_jf|^2)^{1/2}\|_p\le C\|g\|_p\)。
低频块同样由一个 Schwartz 卷积核控制。

反向利用环带上等于 \(1\) 的伴随截断 \(\widetilde\varphi_j\) 及分解恒等式。
与测试函数 \(h\in\mathcal S\) 配对，把 \(J^sf\) 写成块的重构，每项中的 \(J^s\) 除以 \(2^{js}\) 后又是一个一致零阶环带符号。逐点 Cauchy–Schwarz 与\chapref{b2:ch:19}相同的向量值界得到
\[
 |\langle J^sf,h\rangle|
 \le C\left\|\left(|P_0f|^2+\sum2^{2js}|\Delta_jf|^2\right)^{1/2}\right\|_p
          \|h\|_{p'}.
\]
先对有限和证明，再由平方函数尾部的支配收敛及分布重构取极限；对偶性给出反向界。这里没有用习题中的未证明等价，也没有把一般 \(F^s_{p,q}\) 的理论混进来。
\end{proof}
当 \(s\) 为非负整数，导数符号除以 \(J^s\) 满足 Mihlin 条件。反向则以有限个局部符号分解或椭圆符号恒等式控制 \(J^s\)，得到通常 Sobolev 导数范数的等价。仅在 \(p=2\) 时，这种等价可以完全不经过乘子理论而直接由 Plancherel 得到。

% 扩写定位：b2:subsec:220302
整数阶的识别可以把所需符号明确写出。设 \(k\geq1\) 为整数，\(1<p<\infty\)。正向符号
\[
 \frac{(2\pi\mathrm i\xi)^\alpha}{(1+4\pi^2|\xi|^2)^{k/2}},
 \qquad |\alpha|\leq k,
\]
满足 Mihlin 条件，所以 \(\sum_{|\alpha|\leq k}\|\partial^\alpha f\|_p\leq C\|J^kf\|_p\)。
反向先在 \(|\xi|\leq2\) 取一个低频截断，再将 \(|\xi|\geq1\) 的球面分成有限个坐标锥：每片有某个 \(|\xi_j|\geq c|\xi|\)。以光滑齐次分割 \(\chi_j\) 写成
\[
 (1-\chi_0(\xi))(1+4\pi^2|\xi|^2)^{k/2}
 =\sum_{j=1}^d a_j(\xi)(2\pi\mathrm i\xi_j)^k,
\]
其中 \(a_j=(1-\chi_0)\chi_j(1+4\pi^2|\xi|^2)^{k/2}/(2\pi\mathrm i\xi_j)^k\)，在其支集附近定义并平滑延拓为零。分母在该锥上有统一下界，每个 \(a_j\) 满足零阶导数条件。低频符号 \(\chi_0(1+4\pi^2|\xi|^2)^{k/2}\) 的核为 Schwartz 函数。Mihlin 与 Young 给出
\[
 \|J^kf\|_p\leq C\left(\|f\|_p+\sum_{j=1}^d\|\partial_j^kf\|_p\right).
\]
初始等式在 Schwartz 类上成立，再通过分布与光滑逼近得到 \(H^{k,p}=W^{k,p}\)。这里连纯坐标的最高阶导数就已足够控制全部混合导数；它是乘子分解的结论，而不是可以在任意区域或任意边界条件下直接套用的代数事实。

\ProseParagraphBreak

分数阶算子还具有一个反映非局部性的实空间公式。对 \(0<s<2\)，设
\[
 A_{d,s}=\int_{\mathbb R^d}\frac{1-\cos(2\pi z_1)}{|z|^{d+s}}\,\mathrm dz,
 \qquad c_{d,s}=\frac{(2\pi)^s}{A_{d,s}}.
\]
原点处 \(1-\cos(2\pi z_1)\leq C|z|^2\)，无穷处它有界，故 \(0<A_{d,s}<\infty\)；这同时解释范围 \(0<s<2\)。
\begin{proposition}\label{b2:prop:22-fractional-kernel}
对 \(f\in\mathcal S\)，
\[
 (-\Delta)^{s/2}f(x)=\frac{c_{d,s}}2
 \int_{\mathbb R^d}\frac{2f(x)-f(x+y)-f(x-y)}{|y|^{d+s}}\,\mathrm dy.
\]
并且
\[
 \langle(-\Delta)^{s/2}f,f\rangle
 =\frac{c_{d,s}}2
 \iint\frac{|f(x)-f(y)|^2}{|x-y|^{d+s}}\,\mathrm dx\,\mathrm dy.
\]
\end{proposition}
\begin{proof}
对称二阶差在原点是 \(O(|y|^2)\)，远场由函数有界和核尾部可积控制，故第一式绝对收敛。先在 \(\epsilon<|y|<R\) 上作 Fourier 变换，得到乘子
\[
 c_{d,s}\int_{\epsilon<|y|<R}
 (1-\cos(2\pi y\cdot\xi))|y|^{-d-s}\,\mathrm dy.
\]
旋转及换元 \(z=|\xi|y\) 把全积分变成
\(c_{d,s}|\xi|^s A_{d,s}=(2\pi|\xi|)^s\)。可用非负积分的单调收敛，或用 Schwartz 配对的支配收敛取极限，证明第一式。第二式由交换 \(x,y\)、展开平方得到，等价地也可由 Plancherel 与同一个非负乘子积分推导；非负性允许 Tonelli。
\end{proof}
若 \(f\geq0\) 且紧支集，在支集外有
\((-\Delta)^{s/2}f(x)=-c_{d,s}\int f(y)|x-y|^{-d-s}\,\mathrm dy<0\)，除非 \(f=0\)。
因此这个导数一般不保持支集，也不保持非负性；它的二次型却非负。这三个性质并不矛盾：非负的是总能量，而不是每一点的算子输出。与整数 Laplace 算子的局部性相比，这个尾部正是分数阶模型的本质区别。

频率定义中的 Bessel 势也有一个不振荡的实空间核。设 \(s>0\)，定义
\[
 G_s(x)=\frac1{\Gamma(s/2)}
 \int_0^\infty t^{s/2-1}e^{-t}
 (4\pi t)^{-d/2}e^{-|x|^2/(4t)}\,\mathrm dt.
\]
这是把热核作正权平均，权的总质量为一。它与齐次 Riesz 势的区别恰在额外因子 \(e^{-t}\)：大时间被截去，因而不会保留齐次核的长距离幂尾。
\begin{proposition}\label{b2:prop:22-bessel-kernel}
有 \(G_s\geq0\)、\(\int G_s=1\)、\(\widehat G_s(\xi)=(1+4\pi^2|\xi|^2)^{-s/2}\)。故 \(J^{-s}f=G_s*f\)，并在全部 \(1\leq p\leq\infty\) 上满足 \(\|J^{-s}f\|_p\leq\|f\|_p\)。
当 \(0<s<d\) 时，\(G_s(x)\asymp |x|^{s-d}\) 于 \(0<|x|\leq1\)；当 \(s=d\) 时量级为 \(1+\log(1/|x|)\)；当 \(s>d\) 时在原点附近有界。无穷远处三种情形均有指数衰减。
\end{proposition}
\begin{proof}
非负性与 Tonelli 给出总质量一。热核 Fourier 变换为 \(e^{-4\pi^2t|\xi|^2}\)，对 \(t\) 的 gamma 积分得到所述符号，Young 即给出压缩界。

记 \(r=|x|\)。在 \(0<r\leq1\) 时，积分的 \(t\geq1\) 部分有统一有限上界。对 \(0<t<1\)，用 \(e^{-1}\leq e^{-t}\leq1\)，再换元 \(u=r^2/(4t)\)，得到与
\[
 r^{s-d}\int_{r^2/4}^\infty
 u^{(d-s)/2-1}e^{-u}\,\mathrm du
\]
可比的表达。若 \(s<d\)，积分趋于一个正的 gamma 常数；若 \(s=d\)，零点附近是 \(\int_{r^2/4}^1\mathrm du/u\)；若 \(s>d\)，它的下限奇性恰好抵消外面的 \(r^{s-d}\)，结果有界。下界可以分别在固定 \(u\) 区间或 \(r^2/4<u<1\) 中取出，故这些不仅是上界。

当 \(r\geq1\)，在指数中使用 \(t+r^2/(4t)\geq r\)，抽出 \(e^{-r/2}\)。剩下的指数 \(e^{-t/2-r^2/(8t)}\) 把 \(t\to0\) 与 \(t\to\infty\) 的两端都控制；将 \(r\geq1\) 代入后，剩余积分有统一上界或至多一个固定多项式因子，这个因子可再吸收到较慢的指数 \(e^{-cr}\) 中。因此取得远场指数衰减。
\end{proof}
这个核解释了非齐次空间的低频项：常数频率附近的符号有界，所以 \(J^{-s}\) 保持同一 \(L^p\)；齐次符号 \(|\xi|^{-s}\) 在零频发散，改善可积性时必须改变目标指数或规定低频归一化。两者在高频上具有相同阶数，却不能在全空间上互相替代。

\ProseParagraphBreak

核估计还给出 Sobolev 嵌入的另一种表述。若 \(1<p<\infty\)、\(0<s<d/p\)，置 \(1/q=1/p-s/d\)。由 \(G_s(x)\leq C|x|^{s-d}\) 和 HLS，
\[
 \|f\|_q=\|G_s*(J^sf)\|_q\leq C\|J^sf\|_p.
\]
若 \(s>d/p\)，则 \(G_s\in L^{p'}\)：原点幂次条件正是 \((d-s)p'<d\)，临界对数在有限指数也可积，远场由指数衰减处理。因此
\(\|f\|_\infty\leq\|G_s\|_{p'}\|J^sf\|_p\)，并得到连续代表。连续性由 \(L^{p'}\) 中的核平移连续性给出，而且紧支集输入的逼近说明该代表在无穷远趋零。

\ProseParagraphBreak

在临界 \(s=d/p\) 时，核不属于 \(L^{p'}\)，不能原样保留无穷端点。取
\(h(x)=|x|^{-s}(\log(e/|x|))^{-\beta}\mathbf1_{\{|x|<1/2\}}\)，其中 \(1/p<\beta\leq1\)。径向积分给出 \(h\in L^p\)；正核下界则使 \(G_s*h(0)\) 的径向积分成为 \(\int_0^{1/2}r^{-1}(\log(e/r))^{-\beta}\,\mathrm dr=\infty\)。对 \(x\to0\)，Fatou 给出输出任意大的邻域值，故它没有本性有界的代表。这个反例不仅指出一个点的发散，也排除了通过修改零测集恢复 \(L^\infty\) 的可能。

\ProseParagraphBreak

最后，椭圆导数估计与这些嵌入处理不同问题。前者从方程右端恢复同指数的最高阶导数，后者把已经拥有的导数换成更高的可积指数或连续性。应用偏微分方程时，常见路线是先证明解处于 \(H^{s,p}\)，再按 \(s\) 与 \(d/p\) 的相对位置使用嵌入。若只看方程的阶数就声称解连续，会遗漏右端的可积指数及维数这两个决定性参数。

\ProseParagraphBreak

齐次导数与非齐次范数之间还有一个常用的精确比较。对 \(s>0\)、\(1<p<\infty\)，在 Schwartz 类上有
\[
 \|J^sf\|_p\asymp\|f\|_p+\|(-\Delta)^{s/2}f\|_p.
\]
正向用两个乘子：\(J^{-s}\) 以及
\((2\pi|\xi|)^s(1+4\pi^2|\xi|^2)^{-s/2}\)。
后者在原点不一定普通光滑，但定义域只要求 \(\xi\ne0\)，其尺度导数依然受控，故满足 Mihlin 条件。反向取低频截断 \(\chi\)，低频 \(J^s\chi\) 的核可积；高频符号
\[
 (1-\chi(\xi))\frac{(1+4\pi^2|\xi|^2)^{s/2}}{(2\pi|\xi|)^s}
\]
是零阶 Mihlin 符号，把它作用于齐次导数即可恢复高频 \(J^sf\)。两部分合并证明比较；连续延拓把 Schwartz 类的闭包识别为相同的图范数空间。

\ProseParagraphBreak

因此非齐次空间不是拒绝齐次导数，而是在其之外补上低阶 \(L^p\) 大小。高频上的两个定义等价，真正差别留在零频附近。使用能量方法时，可以在已说明的空间范围内自由改写这个图范数；使用全空间齐次逆、缩放不变估计或商掉多项式时，却必须保留齐次定义自己的归一化。一个范数比较不能自动解除所有齐次分布的低频歧义。

\subsection{椭圆型估计}
\label{b2:subsec:220303}

设 \(A=(a_{jk})\) 为实对称正定常矩阵，且
\(\nu|\xi|^2\le\xi^{\mathsf T}A\xi\le\Lambda|\xi|^2\)。
对于紧支撑光滑 \(u\)，令
\(Lu=-\sum a_{jk}\partial_j\partial_k u\)。则
\[
 \widehat{\partial_j\partial_k u}
 =-\frac{\xi_j\xi_k}{\xi^{\mathsf T}A\xi}\widehat{Lu},
 \qquad
 \|D^2u\|_p\le C_{d,p,\nu,\Lambda}\|Lu\|_p.
\]
商是零阶齐次符号，其球面导数由椭圆下界控制；Mihlin 定理给出结论。这是先验估计，既不要求 \(L^{-1}f\) 在 \(L^p\) 中存在，也不将低频问题忽略。

\ProseParagraphBreak
若改用 \(1+L\)，符号
\((1+4\pi^2\xi^{\mathsf T}A\xi)^{-1}\) 及其乘上 \(\xi_j\xi_k\) 的版本均满足零阶界，所以
\[
 \|u\|_p+\|D^2u\|_p\le C\|(1+L)u\|_p.
\]
这可以扩展为 \((1+L)^{-1}:L^p\to W^{2,p}\) 的有界解算子，分布方程成立，唯一性在这一空间中可通过温和分布 Fourier 变换验证。

\ProseParagraphBreak
对于变系数方程，不能把 \(a_{jk}(x)\) 当常数移到 Fourier 符号中，因为 Fourier 变换不与空间乘法交换。局部冻结系数可以形成一个近似常系数问题，但需要小振荡条件、交换子估计和覆盖拼接。\chapref{b2:ch:16}的椭圆正则理论与\appref{b2:app:H}的低正则理论分别处理这些额外困难，本章的常系数证明不替代它们。

常系数模型可以进一步解释定义域与低频唯一性的角色。对正定矩阵 \(A\)，符号 \(1+4\pi^2\xi^{\mathsf T}A\xi\) 没有零点，它的逆及各阶导数至多多项式增长或衰减，因而乘法保持 Schwartz 测试空间。若温和分布 \(u\) 满足 \((1+L)u=0\)，把 Fourier 方程乘以该逆立即得到 \(\widehat u=0\)，故唯一性甚至在整个温和分布类中成立。这里有一个真正可逆的低频项。

\ProseParagraphBreak

对没有常数项的 \(L\)，情况不同。任何常数乃至适当的调和多项式都是齐次方程的温和分布解；若以 \(D^2u\) 的先验界工作，这些低阶自由度不会被看到。在 \(L^p(\mathbb R^d)\)、有限 \(p\) 中，非零常数当然不属于空间，但是否存在具有所需衰减的逆解，仍要检查频率原点的奇性。二阶导数估计已经给出，并不意味着 \(L^{-1}\) 自动为同指数有界算子。

\ProseParagraphBreak

在对角例子 \(L=-\sum_j a_j\partial_j^2\)、\(a_j>0\) 中，换元 \(x_j=\sqrt{a_j}\,y_j\) 将它化为标准 \(-\Delta_y\)。各个导数却分别得到 \(a_j^{-1/2}\) 或 \(a_j^{-1}\) 因子，因此当某个 \(a_j\downarrow0\) 时，椭圆正则估计会退化。一般正定矩阵可先正交对角化，正交变换保持 Lebesgue 测度与欧氏导数长度；最小特征值控制的正是这种退化。记录 \(\nu,\Lambda\) 的依赖有实际意义，而不是可删去的形式装饰。

\ProseParagraphBreak

带参数的预解算子把同一估计组织成一个尺度一致的族。对 \(\lambda>0\)，令 \(u=(\lambda+L)^{-1}f\)。以 \(\xi=\sqrt\lambda\,\zeta\) 归一化后，三个符号
\[
 \frac{\lambda}{\lambda+4\pi^2\xi^{\mathsf T}A\xi},\quad
 \frac{\sqrt\lambda\,2\pi\mathrm i\xi_j}{\lambda+4\pi^2\xi^{\mathsf T}A\xi},\quad
 \frac{-4\pi^2\xi_j\xi_k}{\lambda+4\pi^2\xi^{\mathsf T}A\xi}
\]
都具有与 \(\lambda\) 无关的 Mihlin 常数。于是
\[
 \lambda\|u\|_p+\sqrt\lambda\|\nabla u\|_p+\|D^2u\|_p
 \leq C_{d,p,\nu,\Lambda}\|f\|_p.
\]
符号缩放保持同指数算子范数，这是统一性的关键；直接逐项对未归一化符号求导往往会出现不必要的 \(\lambda^{-1}\) 因子。这里导数和零阶项分别以自己的尺度加权，也正好满足微分阶数的平衡。

\ProseParagraphBreak

令 \(\lambda\downarrow0\) 时，最高阶导数界仍统一，而对 \(\|u\|_p\) 的估计允许增长为 \(\lambda^{-1}\)。因此它没有绕过齐次逆算子的低频障碍。若需要在零参数得到实际解，必须另外控制数据的低频或选择齐次空间；有界参数预解族并不自动给出一个同指数的无质量逆。

\subsection{球面乘子 预告}
\label{b2:subsec:220304}

径向符号只依赖 \(|\xi|\)，仍可能在球面上发生跳跃。
光滑环带截断 \(\chi(|\xi|)\) 的核属于 \(L^1\)；球特征函数 \(\mathbf1_{\{|\xi|\le1\}}\) 则不满足任何跨边界的 Mihlin 光滑条件。
把跳跃平滑成宽度 \(\epsilon\) 的过渡层时，第 \(N\) 阶导数通常增长如 \(\epsilon^{-N}\)。对每个固定 \(\epsilon\) 应用定理可得有界性，却得不到 \(\epsilon\downarrow0\) 的统一界。

\ProseParagraphBreak
Bochner–Riesz 平均以 \((1-|\xi|^2)_+^\delta\) 逐渐消减边界跳跃，其所需光滑阶数与 \(p\)、维数有关。Mihlin 定理给出方便的充分条件，但常不是锐利范围。\chapref{b2:ch:24}用径向核的振荡和曲率研究这一问题；\chapref{b2:ch:23}的限制估计控制频率薄壳上的集中，则说明为什么乘子问题会与曲面延拓相遇。

\ProseParagraphBreak
因此“径向”“有界”“紧支撑”三项都不能单独替代乘子条件。对核使用 Young 不等式、对符号使用尺度导数、对薄壳使用曲面几何，是三条不同证明路线。应用时应明确正在使用哪一条，以及它在端点或退化极限中会失去哪一个条件。

平滑乘子还有一个时间演化模型，可与频率硬截断比较。热算子的符号是 \(e^{-4\pi^2t|\xi|^2}\)。其空间核非负，积分为一，因而在全部 \(L^p\) 上压缩；对每个整数 \(k\geq0\)，对 Gaussian 核求导及换元给出
\[
 \|D^ke^{t\Delta}f\|_p\leq C_k t^{-k/2}\|f\|_p,
 \qquad 1\leq p\leq\infty,\quad t>0.
\]
导数核是一个多项式乘 Gaussian，其 \(L^1\) 范数有限；时间缩放产生 \(t^{-k/2}\)，故该证明连两个端点也包括。单纯的开区间 Mihlin 定理已经足以处理 \(1<p<\infty\)，但核方法在这个特定模型中提供了额外端点。

\ProseParagraphBreak

任意实 \(s\geq0\) 时同样有
\[
 \|e^{t\Delta}f\|_{H^{s,p}}\leq C_{s,p}(1+t^{-s/2})\|f\|_p,
 \qquad1<p<\infty.
\]
在 \(0<t\leq1\) 中，把 \(t^{s/2}(1+4\pi^2|\xi|^2)^{s/2}e^{-4\pi^2t|\xi|^2}\) 以 \(\zeta=\sqrt t\,\xi\) 重写，成为
\((t+4\pi^2|\zeta|^2)^{s/2}e^{-4\pi^2|\zeta|^2}\)。
各个尺度导数在原点附近由 \(|\zeta|^{|\alpha|}\) 抵消可能的负幂，在无穷处由 Gaussian 控制，所以 Mihlin 常数统一；\(t\geq1\) 则用半群性将一次固定时间的平滑与剩余压缩复合。这个推导说明分数阶平滑并不要求先把 \(s\) 取整。

\ProseParagraphBreak

当 \(t\downarrow0\)，正阶范数常数发散，与一般 \(L^p\) 数据可能没有任何正阶导数完全相容。对每个固定正时间得到光滑函数，不意味着这些函数在更强空间中一致趋向原数据。连续半群的初值收敛发生在原来给定的空间，提升正则性则付出时间尺度。

\ProseParagraphBreak

把时间倒过来，符号成为 \(e^{4\pi^2t|\xi|^2}\)，甚至失去平方乘子的本性有界性。取 \(\widehat g\) 支于单位球，并令 \(f_N(x)=e^{2\pi\mathrm iNx_1}g(x)\)。输入平方范数不变，而逆热像的平方范数至少为
\(e^{4\pi^2t(N-1)^2}\|g\|_2\)。适当缩小输入振幅，可以使初值在 \(L^2\) 中趋零而逆演化的像趋于无穷。这个明确的高频检验说明正向平滑不能通过形式上的指数取反得到稳定逆过程。

\ProseParagraphBreak

与球面乘子相比，热符号没有一条频率边界上的跳跃，所有尺度的导数也具有统一控制；与齐次分数积分相比，它在零频不发散。良好的端点压缩、正时间平滑和低频稳定是这些特殊性质共同产生的，不应当归给任意一个写成 \(m(\xi)\) 的符号。

\ProseParagraphBreak

频率过渡层还可以用一个明确的算子范数检验。取固定 \(\chi\in C_c^\infty\)，令 \(m_R(\xi)=\chi(\xi/R)\)。其核为 \(R^d\check\chi(Rx)\)，所以全部 \(R>0\) 的 \(L^1\) 核范数都等于 \(\|\check\chi\|_1\)，在每个 \(L^p\) 上有同一个界。若 \(\chi(0)=1\)，在 \(R\to\infty\) 时，有限指数下这族算子强趋向恒等：先对 Schwartz 输入用核的近似恒等性，再用一致范数和稠密性推广。将频率支撑推向无穷本身并不造成常数发散；关键是过渡轮廓经过尺度归一化后仍固定。

\ProseParagraphBreak

若改用 \(m_\epsilon(\xi)=\theta((|\xi|-1)/\epsilon)\)，其中 \(\theta\) 固定而过渡层紧贴单位球面，归一化整个频率空间不能同时固定球的半径和层的厚度。径向一阶导数在该层上达到 \(\epsilon^{-1}\)，而 \(|\xi|\) 仍约为一，尺度导数条件因此真实地增长。两个模型都叫作光滑截断，却有不同的无量纲参数：前者只有整体半径，后者还有厚度与半径之比。将其混同，会把 Mihlin 常数的一致性误判成一个无条件事实。

\ProseParagraphBreak

在平方空间中只需符号的本性上确界，故上述有界过渡层即便变得越来越薄，仍可具有统一平方界。对其他指数则必须利用核、导数或特殊几何。这个差别再次说明，平方乘子的全部刻画并没有自动延伸到任意 \(L^p\)。讨论逼近一个不连续乘子时，要分别检查固定参数有界、一致参数有界和极限算子的有界性；只有前两项共同存在，才可能通过稠密输入上的极限传递到后一项。


\begin{chaptersummary}
零阶符号的 \(L^2\) 起点由 Plancherel 提供，环带光滑性使核差分可控，Calderón–Zygmund 分解和插值再推广到 \(1<p<\infty\)。Hörmander 条件展示了证明真正需要的积分控制。低频奇性、高频阶数与符号边界几何各有独立作用，不能用一个笼统的“符号光滑”代替核验。
\end{chaptersummary}

\BookExercises

% 自拟：围绕本章已经建立的工具，改变条件或组合方法。
\begin{exercise}\label{b2:ex:22-new-1}
证明非零常数函数 \(m\) 的乘子范数为 \(|m|\)，并说明其核为何不属于普通函数意义下的 \(L^1\)。
\end{exercise}
\begin{solution}
算子为 \(mf\)，故每个 \(L^p\) 上范数均为 \(|m|\)。分布核为 \(m\delta_0\)。若它由 \(L^1\) 函数表示，其 Fourier 变换按 Riemann–Lebesgue 引理应在无穷远趋零，与非零常数矛盾。
\end{solution}

% 自拟：围绕本章已经建立的工具，改变条件或组合方法。
\begin{exercise}\label{b2:ex:22-new-2}
验证 \(m(\xi)=e^{\mathrm i\log|\xi|}\) 满足任意有限阶 Mihlin 条件。它能否连续延拓到原点？
\end{exercise}
\begin{solution}
该符号的模为 \(1\)。每次求导产生一个 \(|\xi|^{-1}\) 因子及球面上的光滑有界函数，归纳得 \(|\partial^\alpha m|\le C_\alpha|\xi|^{-|\alpha|}\)。沿径向趋零时相位不断旋转，无极限；这不影响 \(1<p<\infty\) 的乘子结论。
\end{solution}

% 自拟：围绕本章已经建立的工具，改变条件或组合方法。
\begin{exercise}\label{b2:ex:22-new-3}
证明符号缩放 \(m(r\xi)\) 保持 \(M_p\) 范数，并准确写出共轭公式。
\end{exercise}
\begin{solution}
对 \(D_rf(x)=f(rx)\)，有 \(\widehat{D_rf}(\xi)=r^{-d}\widehat f(\xi/r)\)。直接计算得 \(D_{1/r}T_mD_r=T_{m(r\,\cdot)}\)。两缩放算子的范数为 \(r^{-d/p}\) 与 \(r^{d/p}\)，给出一个方向；用 \(1/r\) 再做一次得到反向，故范数相等。
\end{solution}

% 自拟：围绕本章已经建立的工具，改变条件或组合方法。
\begin{exercise}\label{b2:ex:22-new-4}
设 \(m\) 支撑在固定环带，具有 \(N>d\) 阶连续导数。证明 \(T_m\) 在包括端点的所有 \(L^p\) 上有界。
\end{exercise}
\begin{solution}
分部积分给出 \(|\check m(x)|\le C(1+|x|)^{-N}\)，近原点直接由 \(\|m\|_1\) 控制。由于 \(N>d\)，核可积；Young 不等式给出 \(1\le p\le\infty\) 的界。这个证明只针对单个紧频率环带，并未对无限尺度求和。
\end{solution}

% 自拟：围绕本章已经建立的工具，改变条件或组合方法。
\begin{exercise}\label{b2:ex:22-new-5}
设 \(0<\alpha<d\)。说明为什么 \((2\pi|\xi|)^{-\alpha}\) 不是 \(L^2\) 乘子，却仍可定义分数积分。
\end{exercise}
\begin{solution}
符号在每个零点邻域本质无界，\cref{b2:prop:22-l2}排除 \(L^2\) 有界性。其逆变换是\chapref{b2:ch:19}的 Riesz 位势核，可以在 \(1<p<d/\alpha\) 时定义为 \(L^p\to L^q\)、\(1/q=1/p-\alpha/d\) 的有界算子。空间指数改变，所以两结论一致。
\end{solution}

% 自拟：围绕本章已经建立的工具，改变条件或组合方法。
\begin{exercise}\label{b2:ex:22-new-6}
对 \((1-\Delta)u=f\)，用乘子证明 \(\|u\|_p+\sum_{j,k}\|\partial_j\partial_ku\|_p\le C\|f\|_p\)。
\end{exercise}
\begin{solution}
解符号为 \((1+4\pi^2|\xi|^2)^{-1}\)，二阶导数符号为 \(-4\pi^2\xi_j\xi_k/(1+4\pi^2|\xi|^2)\)。其所有所需尺度导数由有理函数求导直接控制，原点附近有界，高频为零阶或负阶。分别应用 Mihlin 定理并对有限个导数相加即得。
\end{solution}

% 自拟：围绕本章已经建立的工具，改变条件或组合方法。
\begin{exercise}\label{b2:ex:22-new-7}
若 \(\eta\) 为固定光滑环带截断，解释为何 Hörmander 证明同时控制 \(\|h_j\|_1\) 与 \(\|\nabla h_j\|_1\)，而没有假定 \(m\) 可微。
\end{exercise}
\begin{solution}
\(H^s\) 控制转化为 \((1+|x|)^sh_j\) 的 \(L^2\) 控制。梯度对应归一化频率中的坐标乘法；该坐标先乘一个紧支撑光滑截断，在 \(H^s\) 上有界，所以梯度具有相同带权 \(L^2\) 界。因 \(s>d/2\)，逆权重属于 \(L^2\)，Cauchy–Schwarz 给出两项 \(L^1\) 界。它使用频率乘法，不使用 \(m\) 的逐点导数。
\end{solution}

% 自拟：围绕本章已经建立的工具，改变条件或组合方法。
\begin{exercise}\label{b2:ex:22-new-8}
证明一维 \(\mathbf1_{[a,b]}\)、\(a<b\)，是所有 \(1<p<\infty\) 上的乘子，且范数可独立于区间选择。
\end{exercise}
\begin{solution}
Hilbert 变换符号为 \(-\mathrm i\operatorname{sgn}\xi\)，故正半轴投影为 \((I+\mathrm iH)/2\)。调制将半轴边界移到 \(a\) 或 \(b\)，这些操作保持范数。两个相移半轴投影之差的符号为 \(\mathbf1_{[a,b]}\)（端点值无关），算子范数至多 \(1+\|H\|_{p\to p}\)，独立于区间。
\end{solution}
